\section{논의}
\label{sec:discussion}

\subsection{모형 성능 해석}

VAR 모형의 우수한 성능은 거시경제 변수들 간의 구조적 상호의존성을 직접적으로 포착할 수 있기 때문임. GDP, 소비, 투자 등은 단기적으로 순환적 피드백 구조를 형성하며, VAR은 이러한 동시적 상호작용을 벡터 자기회귀 구조로 모델링하므로 단기 예측에서 특히 효과적임.

DFM과 DDFM은 혼합 빈도 데이터 처리 능력을 보유하나, 목표 변수의 데이터 특성에 따라 성능이 크게 달라짐. DFM은 선형 관계가 명확한 경우 효과적이나, 비선형 관계가 중요한 경우 수치적 불안정성을 보임. DDFM의 비선형 인코더는 일부 목표 변수에서 DFM보다 안정적인 성능을 보였으나, 변동성이 큰 변수에서는 한계를 드러냄.

\subsection{모형 선택 가이드라인}

실험 결과를 바탕으로 다음과 같은 모형 선택 가이드라인을 제시함:
\begin{itemize}
    \item \textbf{단기 예측(1-7일)}: VAR 모형이 가장 효과적임
    \item \textbf{중장기 예측(28일 이상)}: ARIMA가 상대적으로 우수한 성능을 보임
    \item \textbf{혼합 빈도 nowcasting}: DFM 또는 DDFM이 적합함. 목표 변수의 데이터 특성에 따라 DFM(선형 관계가 명확한 경우) 또는 DDFM(비선형 관계가 중요한 경우)을 선택해야 함
\end{itemize}

\subsection{연구의 한계점}

본 연구는 다음과 같은 한계점을 가짐: (1) 데이터 품질 및 전처리 과정에서의 정보 손실 가능성, (2) 딥러닝 기반 모형의 해석가능성 제한, (3) 외생 충격 고려 부재, (4) 한국 데이터에 국한된 연구, (5) 일부 조합의 평가 불가능(테스트 세트 크기 부족, 수치적 불안정성).

